% Options for packages loaded elsewhere
% Options for packages loaded elsewhere
\PassOptionsToPackage{unicode}{hyperref}
\PassOptionsToPackage{hyphens}{url}
\PassOptionsToPackage{dvipsnames,svgnames,x11names}{xcolor}
%
\documentclass[
  letterpaper,
]{article}
\usepackage{xcolor}
\usepackage[margin=1in]{geometry}
\usepackage{amsmath,amssymb}
\setcounter{secnumdepth}{5}
\usepackage{iftex}
\ifPDFTeX
  \usepackage[T1]{fontenc}
  % \usepackage[utf8]{inputenc}
  \usepackage{textcomp} % provide euro and other symbols
\else % if luatex or xetex
  \usepackage{unicode-math} % this also loads fontspec
  \defaultfontfeatures{Scale=MatchLowercase}
  \defaultfontfeatures[\rmfamily]{Ligatures=TeX,Scale=1}
\fi
\usepackage{lmodern}
\ifPDFTeX\else
  % xetex/luatex font selection
\fi
% Use upquote if available, for straight quotes in verbatim environments
\IfFileExists{upquote.sty}{\usepackage{upquote}}{}
\IfFileExists{microtype.sty}{% use microtype if available
  \usepackage[]{microtype}
  \UseMicrotypeSet[protrusion]{basicmath} % disable protrusion for tt fonts
}{}
\makeatletter
\@ifundefined{KOMAClassName}{% if non-KOMA class
  \IfFileExists{parskip.sty}{%
    \usepackage{parskip}
  }{% else
    \setlength{\parindent}{0pt}
    \setlength{\parskip}{6pt plus 2pt minus 1pt}}
}{% if KOMA class
  \KOMAoptions{parskip=half}}
\makeatother
% Make \paragraph and \subparagraph free-standing
\makeatletter
\ifx\paragraph\undefined\else
  \let\oldparagraph\paragraph
  \renewcommand{\paragraph}{
    \@ifstar
      \xxxParagraphStar
      \xxxParagraphNoStar
  }
  \newcommand{\xxxParagraphStar}[1]{\oldparagraph*{#1}\mbox{}}
  \newcommand{\xxxParagraphNoStar}[1]{\oldparagraph{#1}\mbox{}}
\fi
\ifx\subparagraph\undefined\else
  \let\oldsubparagraph\subparagraph
  \renewcommand{\subparagraph}{
    \@ifstar
      \xxxSubParagraphStar
      \xxxSubParagraphNoStar
  }
  \newcommand{\xxxSubParagraphStar}[1]{\oldsubparagraph*{#1}\mbox{}}
  \newcommand{\xxxSubParagraphNoStar}[1]{\oldsubparagraph{#1}\mbox{}}
\fi
\makeatother


\usepackage{longtable,booktabs,array}
\newcounter{none} % for unnumbered tables
\usepackage{calc} % for calculating minipage widths
% Correct order of tables after \paragraph or \subparagraph
\usepackage{etoolbox}
\makeatletter
\patchcmd\longtable{\par}{\if@noskipsec\mbox{}\fi\par}{}{}
\makeatother
% Allow footnotes in longtable head/foot
\IfFileExists{footnotehyper.sty}{\usepackage{footnotehyper}}{\usepackage{footnote}}
\makesavenoteenv{longtable}
\usepackage{graphicx}
\makeatletter
\newsavebox\pandoc@box
\newcommand*\pandocbounded[1]{% scales image to fit in text height/width
  \sbox\pandoc@box{#1}%
  \Gscale@div\@tempa{\textheight}{\dimexpr\ht\pandoc@box+\dp\pandoc@box\relax}%
  \Gscale@div\@tempb{\linewidth}{\wd\pandoc@box}%
  \ifdim\@tempb\p@<\@tempa\p@\let\@tempa\@tempb\fi% select the smaller of both
  \ifdim\@tempa\p@<\p@\scalebox{\@tempa}{\usebox\pandoc@box}%
  \else\usebox{\pandoc@box}%
  \fi%
}
% Set default figure placement to htbp
\def\fps@figure{htbp}
\makeatother





\setlength{\emergencystretch}{3em} % prevent overfull lines

\providecommand{\tightlist}{%
  \setlength{\itemsep}{0pt}\setlength{\parskip}{0pt}}



 
\usepackage[]{natbib}
\bibliographystyle{plainnat}


\DeclareUnicodeCharacter{2212}{\ensuremath{-}}
\DeclareUnicodeCharacter{2013}{--}
\DeclareUnicodeCharacter{2014}{---}
\DeclareUnicodeCharacter{2009}{\,}
\DeclareUnicodeCharacter{00A0}{~}
\DeclareUnicodeCharacter{00D7}{\ensuremath{\times}}
\DeclareUnicodeCharacter{2264}{\ensuremath{\leq}}
\DeclareUnicodeCharacter{2265}{\ensuremath{\geq}}
\DeclareUnicodeCharacter{2260}{\ensuremath{\neq}}
\DeclareUnicodeCharacter{2248}{\ensuremath{\approx}}
\DeclareUnicodeCharacter{03B1}{\ensuremath{\alpha}}
\DeclareUnicodeCharacter{03B2}{\ensuremath{\beta}}
\DeclareUnicodeCharacter{03C3}{\ensuremath{\sigma}}
\DeclareUnicodeCharacter{03C7}{\ensuremath{\chi}}
\DeclareUnicodeCharacter{2032}{\ensuremath{'}}

\usepackage{caption}
\captionsetup{labelsep=period}

\setcitestyle{round}

% \AtBeginEnvironment{table}{\footnotesize}
\AtBeginEnvironment{apptbl}{\footnotesize}
% Table font size: \footnotesize by default; table-size.lua emits \nexttablesize{\scriptsize}
% before a table whose div has size="scriptsize", which applies to that table only
\makeatletter
\newcommand\table@defaultsize{\footnotesize}
\let\table@size\table@defaultsize
\newcommand\nexttablesize[1]{\gdef\table@size{#1}}
\AtBeginEnvironment{longtable}{\table@size\global\let\table@size\table@defaultsize}
\makeatother

% Figure notes: note-in-float.lua wraps a figure's Note in \floatnote{...}, which
% prints it after the caption (the float package always sets the caption last)
\makeatletter
\let\floatnote@text\@empty
\newcommand\floatnote[1]{\gdef\floatnote@text{#1}}
\newcommand\fs@plainnote{\fs@plain
  \def\@fs@post{\ifx\floatnote@text\@empty\else
    \vspace\abovecaptionskip\noindent\floatnote@text\par
    \global\let\floatnote@text\@empty\fi}}
\makeatother

% Pin appendix floats in place so they stay next to their Note paragraphs
% (deferred inside \AtBeginDocument because Quarto's \newfloat comes after this file)
\AtBeginDocument{%
  \floatstyle{plainnote}\restylefloat{appfig}%
  \floatplacement{appfig}{H}\floatplacement{apptbl}{H}}

% Caption label for appendix tables typeset as breakable longtables (breakable-apptbl.lua);
% position=top restores the skip between the caption and the table's top rule
\AtBeginDocument{\providecommand*\apptblname{Table W}}
\captionsetup[apptbl]{position=top}

% ACL-style title block: bold title; one column per author (bold name, regular
% affiliation, small monospace email)
\makeatletter
\renewcommand\@maketitle{%
  \newpage\null\vskip 0.5em
  \begin{center}
    {\Large\bfseries \@title\par}
    \vskip 1.5em
    {\normalsize \@author\par}
  \end{center}
  \par\vskip 1em}
\newcommand\aclauthor[3]{% name, affiliation lines, email
  \begin{tabular}[t]{@{}c@{}}
    \textbf{#1}\\[1pt]
    #2\\[1pt]
    {\small\texttt{#3}}
  \end{tabular}}
\makeatother

% Hardcoded PDF authors. Runs at \begin{document}, so it overrides the plain
% \author{} Quarto generates from the YAML (the YAML authors still set PDF metadata).
\AtBeginDocument{%
  \author{%
    \aclauthor{Davood Wadi}
      {Desautels Faculty of Management\\ McGill University\\ Montreal, Quebec, Canada}
      {davood.wadi@hec.ca}%
    \hspace{3em}%
    \aclauthor{Yu Ma}
      {Desautels Faculty of Management\\ McGill University\\ Montreal, Quebec, Canada}
      {yu.ma@mcgill.ca}%
  }%
}
\makeatletter
\@ifpackageloaded{float}{}{\usepackage{float}}
\floatstyle{plain}
\@ifundefined{c@chapter}{\newfloat{apptbl}{h}{loapptbl}}{\newfloat{apptbl}{h}{loapptbl}[chapter]}
\floatname{apptbl}{Table W}
\floatstyle{plaintop}
\restylefloat{apptbl}
\newcommand*\quartoapptblref[1]{Table \hyperref[#1]{W\ref{#1}}}
\@ifpackageloaded{caption}{}{\usepackage{caption}}
\DeclareCaptionLabelFormat{quartoapptblreflabelformat}{#1#2}
\captionsetup[apptbl]{labelformat=quartoapptblreflabelformat}
\newcommand*\listofapptbls{\listof{apptbl}{List of Table Ws}}
\makeatother
\makeatletter
\@ifpackageloaded{float}{}{\usepackage{float}}
\floatstyle{plain}
\@ifundefined{c@chapter}{\newfloat{appfig}{h}{loappfig}}{\newfloat{appfig}{h}{loappfig}[chapter]}
\floatname{appfig}{Figure W}
\newcommand*\quartoappfigref[1]{Figure \hyperref[#1]{W\ref{#1}}}
\@ifpackageloaded{caption}{}{\usepackage{caption}}
\DeclareCaptionLabelFormat{quartoappfigreflabelformat}{#1#2}
\captionsetup[appfig]{labelformat=quartoappfigreflabelformat}
\newcommand*\listofappfigs{\listof{appfig}{List of Figure Ws}}
\makeatother
\makeatletter
\@ifpackageloaded{float}{}{\usepackage{float}}
\floatstyle{plain}
\@ifundefined{c@chapter}{\newfloat{appenv}{h}{loapp}}{\newfloat{appenv}{h}{loapp}[chapter]}
\floatname{appenv}{Web Appendix}
\floatstyle{plaintop}
\restylefloat{appenv}
\newcommand*\quartoappref[1]{Web \hyperref[#1]{Appendix\ref{#1}}}
\@ifpackageloaded{caption}{}{\usepackage{caption}}
\DeclareCaptionLabelFormat{quartoappreflabelformat}{#1#2}
\captionsetup[appenv]{labelformat=quartoappreflabelformat}
\newcommand*\listofappenvs{\listof{appenv}{List of Web Appendixs}}
\makeatother
\makeatletter
\@ifpackageloaded{caption}{}{\usepackage{caption}}
\AtBeginDocument{%
\ifdefined\contentsname
  \renewcommand*\contentsname{Table of contents}
\else
  \newcommand\contentsname{Table of contents}
\fi
\ifdefined\listfigurename
  \renewcommand*\listfigurename{List of Figures}
\else
  \newcommand\listfigurename{List of Figures}
\fi
\ifdefined\listtablename
  \renewcommand*\listtablename{List of Tables}
\else
  \newcommand\listtablename{List of Tables}
\fi
\ifdefined\figurename
  \renewcommand*\figurename{Figure}
\else
  \newcommand\figurename{Figure}
\fi
\ifdefined\tablename
  \renewcommand*\tablename{Table}
\else
  \newcommand\tablename{Table}
\fi
}
\@ifpackageloaded{float}{}{\usepackage{float}}
\floatstyle{ruled}
\@ifundefined{c@chapter}{\newfloat{codelisting}{h}{lop}}{\newfloat{codelisting}{h}{lop}[chapter]}
\floatname{codelisting}{Listing}
\newcommand*\listoflistings{\listof{codelisting}{List of Listings}}
\makeatother
\makeatletter
\makeatother
\makeatletter
\@ifpackageloaded{caption}{}{\usepackage{caption}}
\@ifpackageloaded{subcaption}{}{\usepackage{subcaption}}
\makeatother
\usepackage{bookmark}
\IfFileExists{xurl.sty}{\usepackage{xurl}}{} % add URL line breaks if available
\urlstyle{same}
% fallback for those not using the hyperref driver hyperxmp:
\makeatletter
\@ifundefined{xmpquote}{\newcommand{\xmpquote}[1]{#1}}{}
\makeatother
\hypersetup{
  pdftitle={Can Decision Models Sample Human Behavior? Evaluating Distributional Alignment on Human Survey Populations},
  pdfauthor={Davood Wadi; Yu Ma},
  pdfkeywords={\xmpquote{Decision Models}, \xmpquote{Human Behavior
Simulation}, \xmpquote{Silicon Sampling}, \xmpquote{Ordinal
Psychometrics}, \xmpquote{Distributional ANOVA}},
  colorlinks=true,
  linkcolor={blue},
  filecolor={Maroon},
  citecolor={Blue},
  urlcolor={Blue},
  pdfcreator={LaTeX via pandoc}}


\title{Can Decision Models Sample Human Behavior? Evaluating
Distributional Alignment on Human Survey Populations}
\author{Davood Wadi \and Yu Ma}
\date{2026-10-18}
\begin{document}
\maketitle
\begin{abstract}
Using Large Language Models (LLMs) to simulate human survey responses
(``silicon sampling'') typically relies on generative text sampling or
poorly calibrated token logits, often resulting in mode collapse, prompt
brittleness, and distorted population variance. Recently, \emph{decision
models} (``System One'' architectures) have emerged to map a context
state and typed questions directly to calibrated probability
distributions over answer options in a single prefill forward pass. In
this paper, we investigate whether decision models---and LLMs adapted
via constrained option-letter readouts---can accurately approximate
empirical human response distributions across diverse demographic
subpopulations. Training on respondent-level cross-national survey data
from the International Social Survey Programme (ISSP 2015--2024;
\(N = 458,753\) respondents across 310 country-years) and evaluating
out-of-distribution on US subpopulations from the General Social Survey
(SubPOP; 2,926 subpopulation--question targets), we benchmark
off-the-shelf and fine-tuned decision models against standard causal
LMs. To rigorously evaluate ordinal (Likert) probability mass functions
(PMFs) beyond point estimates and order-invariant Shannon entropy, we
apply an exact analytical framework combining Constraint adherence,
ordinal Consensus/Dissension, 1-Wasserstein (Earth Mover's) distance,
and Distributional ANOVA across demographic factors. We show where
decision models succeed in capturing human marginal distributions and
inter-group contrasts, and where both architectures systematically fail
to reflect human behavioral heterogeneity.
\end{abstract}


\section{Introduction}\label{sec-intro}

Simulating human attitudes, preferences, and survey responses with
foundation models---often termed \emph{silicon sampling} or \emph{in
silico} surveying---has attracted widespread interest across
computational social science, economics, and behavioral research
\citep{argyle2023out, horton2023large, brand2023using, sarstedt2024using, pmlr-v202-santurkar23a, durmus2024towards}.
However, using standard autoregressive Large Language Models (LLMs) as
proxies for human populations faces three fundamental methodological
bottlenecks:

\begin{enumerate}
\def\labelenumi{\arabic{enumi}.}
\tightlist
\item
  \textbf{The Generative Sampling Bottleneck:} Autoregressive text
  generation requires repeated Monte Carlo draws to estimate a
  subpopulation's response distribution \citep{wadi2025monte},
  introducing sampling noise, sensitivity to decoding temperature, and
  ``fence-sitting'' refusals
  \citep{bhatia2025value, liu2025badworktimecrosscultural}.
\item
  \textbf{The Calibration and Homogenization Bottleneck:} Next-token
  probabilities in instruction-tuned LLMs are often miscalibrated
  (typically overconfident or collapsed onto a single mode) and
  frequently fail to reflect true inter-group variance across
  demographic attributes such as gender, political ideology, or religion
  \citep{pmlr-v202-santurkar23a, deas-mckeown-2025-artificial}.
\item
  \textbf{The Ordinal Evaluation Bottleneck:} Standard NLP metrics
  (e.g., accuracy of the argmax choice, Shannon entropy
  \citep{farquhar2024detecting}, or KL divergence alone) treat survey
  options as unordered categories. On Likert scales
  \citep{batterton2017likert}, Shannon entropy cannot distinguish
  between a population polarized at the extremes (\emph{Strongly
  Disagree} vs.~\emph{Strongly Agree}) and one concentrated around
  adjacent moderate options (\emph{Neutral} vs.~\emph{Agree})
  \citep{wadi2026every}.
\end{enumerate}

Recently, a new class of architectures---\textbf{decision models} (or
``System One'' models, drawing on dual-process theory
\citep{kahneman2011thinking, evans2013dual}, such as TypeSafe's Jev,
Intern-Decision, and Kev
\citep{deusser2026evaluatingbenchmarkingmodeljev, li2026llmasjevllmsjevstyledecision, geuter2026releasinginnerjevsimple})---has
been introduced to replace text generation with structured, single-pass
probabilistic decision-making. Given a state (e.g., a respondent
persona) and a set of typed questions (\emph{score}, \emph{choice},
\emph{noul}), a decision model executes a single prefill-only forward
pass and outputs a probability mass function (PMF) over the supplied
answer options via a dedicated pointer head or restricted option-token
readout.

Because decision models are explicitly architected to output calibrated
distributions over discrete options without text generation, a natural
question arises: \textbf{Can decision models be used to sample human
behavior?} Specifically, how well do off-the-shelf and domain-adapted
decision models---compared to standard LLMs
\citep{qwen3.5, suh-etal-2025-language}---approximate empirical human
response distributions and demographic contrasts?

\subsection{Contributions}\label{contributions}

\begin{enumerate}
\def\labelenumi{\arabic{enumi}.}
\tightlist
\item
  \textbf{Empirical Benchmark of Decision Models vs.~LLMs on Human
  Populations:} We evaluate pointer-head decision models (Kev-0.8B,
  Kev-4B, Kev-9B) and next-token letter-readout LLMs (Qwen3.5
  \citep{qwen3.5}) across zero-shot, LoRA, and full fine-tuning regimes,
  training on 10 waves of respondent-level ISSP data (\(458,753\)
  individuals) and testing out-of-distribution on US subpopulations from
  SubPOP \citep{suh-etal-2025-language}.
\item
  \textbf{Exact Ordinal Psychometric Evaluation Framework:} Building on
  our analytical factorial framework for ordinal distributions
  \citep{wadi2026every}, we evaluate model-predicted PMFs against
  observed human subpopulation shares across four complementary
  dimensions: (i) \emph{Constraint} (valid option mass
  \citep{cronbach1950further}), (ii) \emph{Distributional Fidelity} (KL
  divergence, Jensen--Shannon divergence, and normalized 1-Wasserstein
  distance \citep{villani2008optimal}), (iii) \emph{Ordinal Dispersion}
  (Consensus and Dissension \citep{tastle2007consensus}), and (iv)
  \emph{Construct \& Contrast Decomposition} (Distributional ANOVA
  isolating demographic main effects and inter-group gaps
  \citep{wadi2026interplay, wadi2026every}).
\item
  \textbf{Diagnosis of Failure Modes in Human Simulation:} We
  characterize when models suffer from ordinal scale compression
  \citep{gao2026choicesfewerdecisionsordinalscale} or collapse to
  marginal national averages (ignoring persona conditioning) versus when
  fine-tuning on respondent-level attribute subsets recovers genuine
  human demographic divergences.
\end{enumerate}

\begin{center}\rule{0.5\linewidth}{0.5pt}\end{center}

\section{Background and Related Work}\label{sec-related}

\subsection{Simulating Human Populations with LLMs (``Silicon
Sampling'')}\label{simulating-human-populations-with-llms-silicon-sampling}

\begin{itemize}
\tightlist
\item
  Early work conditioning LLMs on demographic backstories to approximate
  human economic, consumer, and public opinion surveys
  \citep{argyle2023out, horton2023large, brand2023using, sarstedt2024using, goli2024frontiers, pmlr-v202-santurkar23a, durmus2024towards}.
\item
  Fine-tuning and distributional alignment on human survey distributions
  \citep{suh-etal-2025-language, huang-etal-2026-distribution, ku2026siliconsamplingcrosssurveytransfer}.
\item
  Cognitive and behavioral biases in LLMs compared to human
  decision-makers
  \citep{binz2023using, hagendorff2023human, macmillan2024ir, cheung2025large, echterhoff2024cognitive, mina2025cognitive}.
\item
  Known failure modes: demographic stereotyping, flattening of
  within-group variance, sensitivity to prompt formatting and option
  order
  \citep{shu-etal-2024-dont, kabir-etal-2025-break, koo2024benchmarking},
  low ecological validity of generative outputs
  \citep{jung-etal-2026-psychometric, deas-mckeown-2025-artificial}, and
  poor alignment with non-WEIRD populations
  \citep{henrich2010weirdest, zhao-etal-2024-worldvaluesbench}.
\end{itemize}

\subsection{From Generative Sampling to Exact Token Probabilities and
Decision
Models}\label{from-generative-sampling-to-exact-token-probabilities-and-decision-models}

\begin{itemize}
\tightlist
\item
  Extracting exact next-token probabilities over answer options rather
  than sampling free text
  \citep{kuribayashi-etal-2024-psychometric, bentegeac2025token, kadavath2022language, wadi2026every}.
\item
  \textbf{Decision Models (``System One''):} Pretrained causal LM
  backbones \citep{qwen3.5} adapted to ingest a context state and typed
  questions (\emph{noul}, \emph{choice}, \emph{score}) and produce
  probabilities in one prefill pass
  \citep{deusser2026evaluatingbenchmarkingmodeljev, asseri2026calibratedwhompersonalanguage, li2026kitescalingjevpopulation}:

  \begin{itemize}
  \tightlist
  \item
    \emph{Readout Style A (Restricted Next-Token Logits):} Mapping
    options to single-token symbols (``A'', ``B'', \(\dots\)) and taking
    a softmax over allowed option logits at the answer position (e.g.,
    Intern-Decision, JevK5, and option-letter LLM readouts
    \citep{li2026llmasjevllmsjevstyledecision, geuter2026releasinginnerjevsimple}).
  \item
    \emph{Readout Style B (Pointer Head):} Scoring the hidden state at
    each option's closing token against the question's final token via
    learned linear projections (e.g., Kev), decoupling option scoring
    from vocabulary token priors.
  \end{itemize}
\item
  \textbf{Ordinal Scale Compression in Decision Models:} Recent evidence
  that direct-decision models systematically underuse ordinal rating
  scales \citep{gao2026choicesfewerdecisionsordinalscale}, motivating
  rigorous ordinal distributional evaluation.
\end{itemize}

\subsection{Evaluating Ordinal Probability
Distributions}\label{evaluating-ordinal-probability-distributions}

\begin{itemize}
\tightlist
\item
  Why point estimates (expected values) and categorical metrics (Shannon
  entropy \citep{farquhar2024detecting}, top-1 accuracy) fail on Likert
  scales \citep{batterton2017likert, joshi2015likert, wadi2026every}.
\item
  Ordinal dispersion via Consensus and Dissension
  \citep{tastle2007consensus} and optimal transport (1-Wasserstein
  distance \citep{villani2008optimal}) on discrete ordered supports.
\end{itemize}

\begin{center}\rule{0.5\linewidth}{0.5pt}\end{center}

\section{Problem Formulation and Architectures}\label{sec-problem}

\subsection{Task Definition: Subpopulation Distributional
Prediction}\label{task-definition-subpopulation-distributional-prediction}

Let a human survey dataset consist of respondents indexed by
\(i \in \{1, \dots, N\}\). Each respondent has a vector of demographic
and background attributes \(\mathbf{a}_i\) (e.g., country, sex, age,
education, religion, political party) and answers to survey questions
\(q \in \mathcal{Q}\). Each question \(q\) has \(K_q\) discrete response
options \(\mathcal{Y}_q = \{y_1, \dots, y_{K_q}\}\). For ordinal Likert
(\emph{score}) items, \((\mathcal{Y}_q, \preceq)\) is totally ordered
from lowest to highest intensity.

A \textbf{persona} (or experimental condition) \(\boldsymbol{\lambda}\)
specifies a subset of attribute--value pairs (e.g., ``Country: United
States; Political party: Republican''). Let
\(\mathcal{I}(\boldsymbol{\lambda}, q)\) denote the set of human
respondents matching persona \(\boldsymbol{\lambda}\) who provided a
valid answer to question \(q\), with group size
\(n_{\boldsymbol{\lambda}, q} = |\mathcal{I}(\boldsymbol{\lambda}, q)|\).
The empirical human response distribution (target PMF) is:
\begin{equation}\protect\phantomsection\label{eq-human-pmf}{
P_{\mathrm{human}}(Y_{q, \boldsymbol{\lambda}} = y_k) = \frac{1}{n_{\boldsymbol{\lambda}, q}} \sum_{i \in \mathcal{I}(\boldsymbol{\lambda}, q)} \mathbb{I}(y_{i, q} = y_k), \quad k \in \{1, \dots, K_q\}.
}\end{equation}

Given the compiled state representation of \(\boldsymbol{\lambda}\) and
the question specification \(q\), a model \(M\) outputs logits
\(\mathbf{z}_{q, \boldsymbol{\lambda}} \in \mathbb{R}^{K_q}\) in a
single forward pass, defining the predicted PMF without temperature
rescaling (\(T = 1\)):
\begin{equation}\protect\phantomsection\label{eq-model-pmf}{
P_{\mathrm{model}}(Y_{q, \boldsymbol{\lambda}} = y_k) = \operatorname{softmax}(\mathbf{z}_{q, \boldsymbol{\lambda}})_k = \frac{\exp(z_{q, \boldsymbol{\lambda}, k})}{\sum_{j=1}^{K_q} \exp(z_{q, \boldsymbol{\lambda}, j})}.
}\end{equation}

\subsection{Model Architectures
Compared}\label{model-architectures-compared}

We compare two distinct mechanisms for turning a causal language model
backbone into a distributional survey predictor:

\subsubsection{Pointer-Head Decision Model
(Kev)}\label{pointer-head-decision-model-kev}

In the pointer-head architecture, the persona state and question options
are compiled into a structured prompt. Rather than reading from the
language modeling vocabulary head, a lightweight pointer head projects
the final-layer hidden state at the question's closing token
(\(\mathbf{h}_{\mathrm{q}}\)) and at each option \(k\)'s closing token
(\(\mathbf{h}_{\mathrm{opt}, k}\)):
\begin{equation}\protect\phantomsection\label{eq-pointer-head}{
z_{q, \boldsymbol{\lambda}, k} = \frac{(\mathbf{W}_q \mathbf{h}_{\mathrm{q}})^\top (\mathbf{W}_k \mathbf{h}_{\mathrm{opt}, k})}{\sqrt{d_{\mathrm{head}}}}.
}\end{equation} Because options are scored via their contextualized span
representations rather than fixed option-letter tokens, the pointer head
natively handles arbitrary option counts and avoids token-id frequency
biases.

\subsubsection{Constrained Next-Token Letter
LLM}\label{constrained-next-token-letter-llm}

In the standard LLM readout, the persona and question options (labelled
``A.'', ``B.'', \(\dots\)) are formatted using the model's native chat
template with reasoning disabled, positioning the next predicted token
immediately at the assistant's response. Let \(\mathcal{V}\) be the full
vocabulary and
\(\mathcal{V}_{\mathrm{val}} = \{t_1, \dots, t_{K_q}\} \subset \mathcal{V}\)
be the single-token identifiers corresponding to option letters ``A''
through the \(K_q\)-th letter. The raw vocabulary logits at the final
prompt position yield
\(z_{q, \boldsymbol{\lambda}, k} = \ell_{\mathrm{LM}}(t_k)\).

\subsection{Fine-Tuning on Respondent-Level Human
Data}\label{fine-tuning-on-respondent-level-human-data}

To test whether discrepancies between models and humans stem from
architectural limitations or a lack of behavioral calibration during
pretraining, we fine-tune both model families on respondent-level survey
data using soft cross-entropy loss against observed subpopulation
shares: \begin{equation}\protect\phantomsection\label{eq-soft-ce}{
\mathcal{L}(\theta) = -\sum_{k=1}^{K_q} P_{\mathrm{human}}(Y_{q, \boldsymbol{\lambda}} = y_k) \log P_{\mathrm{model}}(Y_{q, \boldsymbol{\lambda}} = y_k; \theta),
}\end{equation} which is equivalent to minimizing the Kullback--Leibler
divergence
\(D_{\mathrm{KL}}(P_{\mathrm{human}} \,\|\, P_{\mathrm{model}})\) up to
the constant human target entropy \(H(P_{\mathrm{human}})\).

\begin{center}\rule{0.5\linewidth}{0.5pt}\end{center}

\section{Methodology: Analytical Evaluation of Ordinal
Distributions}\label{sec-methodology}

To evaluate how well \(P_{\mathrm{model}}\) approximates
\(P_{\mathrm{human}}\), standard metrics such as top-1 accuracy or
expected value comparisons are insufficient: two distributions can share
the exact same mean while representing opposite sociological realities
(e.g., unimodal consensus vs.~bimodal polarization). Building on our
analytical factorial framework for ordinal scales, we evaluate models
across four layers.

\subsection{Layer 1: Constraint Adherence (Valid Option
Mass)}\label{layer-1-constraint-adherence-valid-option-mass}

For standard LLMs projecting onto the full vocabulary \(\mathcal{V}\),
the model may assign probability mass to invalid tokens (e.g.,
conversational preambles, refusals, or out-of-range symbols). We define
the \textbf{Constraint Failure Rate} for prompt
\(x_{q, \boldsymbol{\lambda}}\) as:
\begin{equation}\protect\phantomsection\label{eq-failure-rate}{
\mathrm{FailureRate}(x_{q, \boldsymbol{\lambda}}) = 1 - \sum_{t \in \mathcal{V}_{\mathrm{val}}} P_{\mathrm{raw}}(t \mid x_{q, \boldsymbol{\lambda}}).
}\end{equation} Pointer-head decision models satisfy
\(\mathrm{FailureRate} \equiv 0\) by construction, whereas LLMs must
learn or be prompted to concentrate mass on
\(\mathcal{V}_{\mathrm{val}}\) before renormalization.

\subsection{Layer 2: Distributional Fidelity (Categorical and Ordinal
Distances)}\label{layer-2-distributional-fidelity-categorical-and-ordinal-distances}

For all question types (\emph{choice}, \emph{score}, \emph{noul}), we
measure information-theoretic divergence in bits:
\begin{equation}\protect\phantomsection\label{eq-kl}{
D_{\mathrm{KL}}(P_{\mathrm{human}} \,\|\, P_{\mathrm{model}}) = \sum_{k=1}^{K_q} P_{\mathrm{human}}(y_k) \log_2 \frac{P_{\mathrm{human}}(y_k)}{P_{\mathrm{model}}(y_k)},
}\end{equation} along with the symmetric, bounded
\textbf{Jensen--Shannon Divergence}
\(D_{\mathrm{JS}}(P_{\mathrm{human}}, P_{\mathrm{model}}) \in [0, 1]\).

However, \(D_{\mathrm{KL}}\) and \(D_{\mathrm{JS}}\) are
permutation-invariant: they penalize a shift from \emph{Strongly Agree}
to \emph{Agree} identically to a shift from \emph{Strongly Agree} to
\emph{Strongly Disagree}. For ordinal Likert (\emph{score}) items with
ordered levels \(y_1 \prec \dots \prec y_{K_q}\) mapped to
\(\{0, 1, \dots, K_q - 1\}\), we compute the \textbf{1-Wasserstein
Distance (Earth Mover's Distance)} via the \(L_1\) distance between
cumulative distribution functions (CDFs) \(F_{\mathrm{human}}\) and
\(F_{\mathrm{model}}\), normalized by the scale range \(K_q - 1\):
\begin{equation}\protect\phantomsection\label{eq-wasserstein}{
W_1(P_{\mathrm{human}}, P_{\mathrm{model}}) = \frac{1}{K_q - 1} \sum_{k=1}^{K_q - 1} \left| \sum_{j=1}^{k} P_{\mathrm{human}}(y_j) - \sum_{j=1}^{k} P_{\mathrm{model}}(y_j) \right| \in [0, 1].
}\end{equation} \(W_1\) measures the minimum fraction of the scale range
across which probability mass must be transported to transform the
model's distribution into the empirical human distribution.

\subsection{Layer 3: Ordinal Consensus and
Dissension}\label{layer-3-ordinal-consensus-and-dissension}

To test whether models replicate the internal agreement or polarization
of a human subpopulation (rather than merely matching the mean or
artificially flattening/peaking the distribution), we apply the ordinal
\textbf{Consensus (\(\mathrm{Cns}\))} metric \citep{tastle2007consensus}
and its multivariate generalization. For a single ordinal item with
support \(\mathcal{Y} = \{1, \dots, K_q\}\), mean
\(\mu = \sum_{k=1}^{K_q} k \, P(y_k)\), and scale width
\(d_{\max} = K_q - 1\):
\begin{equation}\protect\phantomsection\label{eq-consensus-univariate}{
\mathrm{Cns}(Y) = 1 + \sum_{k=1}^{K_q} P(y_k) \log_2 \left( 1 - \frac{|k - \mu|}{d_{\max}} \right), \quad \mathrm{Dsn}(Y) = 1 - \mathrm{Cns}(Y).
}\end{equation} Across a multi-item battery
\(\mathbf{Y}_{\boldsymbol{\lambda}} = (Y_{1, \boldsymbol{\lambda}}, \dots, Y_{M, \boldsymbol{\lambda}}) \in \mathcal{Y}^M\)
with mean vector \(\boldsymbol{\mu}\) and maximum diagonal distance
\(d_{\max} = \sqrt{M}(y_{\max} - y_{\min})\), multivariate Consensus is:
\begin{equation}\protect\phantomsection\label{eq-consensus-multivariate}{
\mathrm{Cns}(\mathbf{Y}_{\boldsymbol{\lambda}}) = 1 + \sum_{\mathbf{y} \in \mathcal{Y}^M} P(\mathbf{y}) \log_2 \left( 1 - \frac{\|\mathbf{y} - \boldsymbol{\mu}\|_2}{d_{\max}} \right).
}\end{equation} Comparing
\(\Delta \mathrm{Cns} = \mathrm{Cns}(P_{\mathrm{model}}) - \mathrm{Cns}(P_{\mathrm{human}})\)
directly reveals whether a model is \textbf{over-concentrated}
(\(\Delta \mathrm{Cns} > 0\), mode collapse) or \textbf{over-dispersed /
polarized} (\(\Delta \mathrm{Cns} < 0\)) relative to real human
respondents.

\subsection{Layer 4: Construct Convolution and Distributional Factorial
Decomposition}\label{layer-4-construct-convolution-and-distributional-factorial-decomposition}

When evaluating how demographic factors (e.g.,
\(\text{Party} \in \{\text{Democrat}, \text{Republican}\}\),
\(\text{Sex} \in \{\text{Male}, \text{Female}\}\)) shift attitudes
across a single item or a convolved multi-item construct
\(P_{S_{\boldsymbol{\lambda}}} = P_{Y_{1, \boldsymbol{\lambda}}} \circledast \cdots \circledast P_{Y_{M, \boldsymbol{\lambda}}}\),
we apply \textbf{Distributional ANOVA} (distributional Hoeffding
decomposition):
\begin{equation}\protect\phantomsection\label{eq-anova-identity}{
\mathbb{E}[S_{\boldsymbol{\lambda}}] = \underbrace{\mathbb{E}[S_0]}_{\text{Grand Mean}} + \sum_{c \in \mathcal{C}} \underbrace{\mathbb{E}[E_c(\lambda_c)]}_{\text{Main Effects}} + \sum_{\substack{U \subseteq \mathcal{C} \\ |U| \ge 2}} \underbrace{\mathbb{E}[E_U(\boldsymbol{\lambda}_U)]}_{\text{Interactions}}.
}\end{equation} For any two subpopulations
\(\boldsymbol{\lambda}_a, \boldsymbol{\lambda}_b\) differing on factor
\(c\) (e.g., Democrat vs.~Republican), we compare the model's predicted
behavioral contrast against the true human contrast using:

\begin{itemize}
\tightlist
\item
  \textbf{Inter-Group Contrast Error (\(\Delta W_1\)):} How well the
  model reproduces the Wasserstein gap
  \(W_1(P_{\mathrm{human}}(\cdot \mid \boldsymbol{\lambda}_a), P_{\mathrm{human}}(\cdot \mid \boldsymbol{\lambda}_b))\)
  between the two human groups.
\item
  \textbf{Effect Summary Metrics:} Mean Shift (\(\mathbb{E}[E]\)),
  Predictive Dispersion (\(\mathrm{SD}(E)\)), Discrete Probability of
  Direction
  (\(\mathrm{dPD} = \max(\mathbb{P}(E > 0), \mathbb{P}(E < 0))\)), and
  Signal-to-Noise Ratio
  (\(\mathrm{SNR} = |\mathbb{E}[E]| / \mathrm{SD}(E)\)).
\end{itemize}

\begin{center}\rule{0.5\linewidth}{0.5pt}\end{center}

\section{Experimental Setup}\label{sec-experiments}

\subsection{Datasets}\label{datasets}

\begin{longtable}[]{@{}
  >{\raggedright\arraybackslash}p{(\linewidth - 8\tabcolsep) * \real{0.2000}}
  >{\raggedright\arraybackslash}p{(\linewidth - 8\tabcolsep) * \real{0.2000}}
  >{\raggedright\arraybackslash}p{(\linewidth - 8\tabcolsep) * \real{0.2000}}
  >{\raggedright\arraybackslash}p{(\linewidth - 8\tabcolsep) * \real{0.2000}}
  >{\raggedright\arraybackslash}p{(\linewidth - 8\tabcolsep) * \real{0.2000}}@{}}
\caption{Summary of human survey datasets used for training,
calibration, and out-of-distribution
evaluation.}\label{tbl-datasets}\tabularnewline
\toprule\noalign{}
\begin{minipage}[b]{\linewidth}\raggedright
Split / Role
\end{minipage} & \begin{minipage}[b]{\linewidth}\raggedright
Dataset
\end{minipage} & \begin{minipage}[b]{\linewidth}\raggedright
Source \& Scope
\end{minipage} & \begin{minipage}[b]{\linewidth}\raggedright
Scale
\end{minipage} & \begin{minipage}[b]{\linewidth}\raggedright
Question Types
\end{minipage} \\
\midrule\noalign{}
\endfirsthead
\toprule\noalign{}
\begin{minipage}[b]{\linewidth}\raggedright
Split / Role
\end{minipage} & \begin{minipage}[b]{\linewidth}\raggedright
Dataset
\end{minipage} & \begin{minipage}[b]{\linewidth}\raggedright
Source \& Scope
\end{minipage} & \begin{minipage}[b]{\linewidth}\raggedright
Scale
\end{minipage} & \begin{minipage}[b]{\linewidth}\raggedright
Question Types
\end{minipage} \\
\midrule\noalign{}
\endhead
\bottomrule\noalign{}
\endlastfoot
\textbf{Training} & ISSP (2015--2024) & 10 thematic modules across 310
country-years & 458,753 respondents & \emph{score}, \emph{choice},
\emph{noul} \\
\textbf{Calibration} & ISSP Held-Out Slots & 1\% of (country, question)
slots per year; country-level \& 1-attribute marginal groups & 204 slots
& \emph{score}, \emph{choice}, \emph{noul} \\
\textbf{Out-of-Distribution Test} & SubPOP (GSS) & US General Social
Survey (GSS), 22 US demographic subpopulations & 133 questions; 2,926
targets & 64 \emph{score} (Likert), 69 \emph{choice} \\
\end{longtable}

\subsubsection{Dynamic Respondent-Level Sampling
(ISSP)}\label{dynamic-respondent-level-sampling-issp}

Rather than aggregating respondents into a fixed small set of
demographic cells (which averages away attributes like religion, party,
or income), our training sampler dynamically constructs
persona--question records on the fly from respondent-level ISSP data:

\begin{enumerate}
\def\labelenumi{\arabic{enumi}.}
\tightlist
\item
  Draw a (year, country, question) slot uniformly at random.
\item
  Draw attribute count \(k \sim \mathrm{Uniform}(0, A_{\text{avail}})\)
  from up to 16 background attributes (sex, age, education, religion,
  religious attendance, party vote, left--right party position, turnout,
  household income, region, urban/rural, marital status, main activity,
  union membership, subjective social status, ancestry).
\item
  Sample an observed attribute combination with at least \(n_{\min}\)
  matching respondents and compute their exact empirical answer shares
  as the training target.
\end{enumerate}

\subsubsection{Out-of-Distribution Test Set (SubPOP /
GSS)}\label{out-of-distribution-test-set-subpop-gss}

We evaluate all models on the SubPOP evaluation split
\citep{suh-etal-2025-language}, comprising 133 US General Social Survey
questions across 22 single-attribute US subpopulations (2,926 total
target distributions; 64 ordinal Likert \emph{score} items and 69
categorical \emph{choice} items). No temperature scaling (\(T=1\)) is
applied at test time.

\subsection{Models and Training
Regimes}\label{models-and-training-regimes}

We evaluate the following model configurations:

\begin{enumerate}
\def\labelenumi{\arabic{enumi}.}
\tightlist
\item
  \textbf{Off-the-Shelf Decision Models:} Kev-0.8B, Kev-4B, and Kev-9B
  (pretrained Qwen3.5 backbones with released LoRA adapters and pointer
  heads trained on general decision benchmarks).
\item
  \textbf{Fine-Tuned Decision Models (Kev):}

  \begin{itemize}
  \tightlist
  \item
    \emph{Adapter / LoRA Fine-Tuning:} Continuing the released adapter
    or training a fresh rank-\(R\) LoRA adapter (\(R \in \{16, 64\}\))
    plus the pointer head on ISSP.
  \item
    \emph{Layer Unfreezing / Full Fine-Tuning:} Unfreezing the top-\(L\)
    transformer layers or all backbone weights alongside the pointer
    head.
  \item
    \emph{Training Scope Ablations:} All 310 country-years vs.~US-only
    respondents; persona complexity (\(k\) attributes) and minimum group
    size (\(n_{\min} \in \{1, 10, 100\}\)).
  \end{itemize}
\item
  \textbf{Constrained Next-Token LLMs:} Qwen3.5-0.8B (and larger
  variants) evaluated zero-shot and fine-tuned end-to-end through its
  vocabulary head over option letters ``A''--``L''.
\end{enumerate}

\begin{center}\rule{0.5\linewidth}{0.5pt}\end{center}

\section{Results}\label{sec-results}

\subsection{Overall Distributional Fidelity on Human Subpopulations
(RQ1)}\label{overall-distributional-fidelity-on-human-subpopulations-rq1}

\emph{TODO: Populate Table~\ref{tbl-main-results} with test-set metrics
from the evaluation runs.}

\begin{longtable}[]{@{}
  >{\raggedright\arraybackslash}p{(\linewidth - 14\tabcolsep) * \real{0.1000}}
  >{\raggedright\arraybackslash}p{(\linewidth - 14\tabcolsep) * \real{0.1000}}
  >{\raggedleft\arraybackslash}p{(\linewidth - 14\tabcolsep) * \real{0.1333}}
  >{\raggedleft\arraybackslash}p{(\linewidth - 14\tabcolsep) * \real{0.1333}}
  >{\raggedleft\arraybackslash}p{(\linewidth - 14\tabcolsep) * \real{0.1333}}
  >{\raggedleft\arraybackslash}p{(\linewidth - 14\tabcolsep) * \real{0.1333}}
  >{\raggedleft\arraybackslash}p{(\linewidth - 14\tabcolsep) * \real{0.1333}}
  >{\raggedleft\arraybackslash}p{(\linewidth - 14\tabcolsep) * \real{0.1333}}@{}}
\caption{Out-of-distribution performance on SubPOP (2,926 subpopulation
targets across 133 GSS questions). KL and JS divergences are in bits;
\(W_1\) is normalized by scale range (\(\%\) of
scale).}\label{tbl-main-results}\tabularnewline
\toprule\noalign{}
\begin{minipage}[b]{\linewidth}\raggedright
Model
\end{minipage} & \begin{minipage}[b]{\linewidth}\raggedright
Regime
\end{minipage} & \begin{minipage}[b]{\linewidth}\raggedleft
All: \(D_{\mathrm{KL}}\) (\(\downarrow\))
\end{minipage} & \begin{minipage}[b]{\linewidth}\raggedleft
All: \(D_{\mathrm{JS}}\) (\(\downarrow\))
\end{minipage} & \begin{minipage}[b]{\linewidth}\raggedleft
Choice: \(D_{\mathrm{KL}}\) (\(\downarrow\))
\end{minipage} & \begin{minipage}[b]{\linewidth}\raggedleft
Score: \(D_{\mathrm{KL}}\) (\(\downarrow\))
\end{minipage} & \begin{minipage}[b]{\linewidth}\raggedleft
Score: \(W_1\) (\(\downarrow\))
\end{minipage} & \begin{minipage}[b]{\linewidth}\raggedleft
\(\Delta \mathrm{Cns}\) (\(\to 0\))
\end{minipage} \\
\midrule\noalign{}
\endfirsthead
\toprule\noalign{}
\begin{minipage}[b]{\linewidth}\raggedright
Model
\end{minipage} & \begin{minipage}[b]{\linewidth}\raggedright
Regime
\end{minipage} & \begin{minipage}[b]{\linewidth}\raggedleft
All: \(D_{\mathrm{KL}}\) (\(\downarrow\))
\end{minipage} & \begin{minipage}[b]{\linewidth}\raggedleft
All: \(D_{\mathrm{JS}}\) (\(\downarrow\))
\end{minipage} & \begin{minipage}[b]{\linewidth}\raggedleft
Choice: \(D_{\mathrm{KL}}\) (\(\downarrow\))
\end{minipage} & \begin{minipage}[b]{\linewidth}\raggedleft
Score: \(D_{\mathrm{KL}}\) (\(\downarrow\))
\end{minipage} & \begin{minipage}[b]{\linewidth}\raggedleft
Score: \(W_1\) (\(\downarrow\))
\end{minipage} & \begin{minipage}[b]{\linewidth}\raggedleft
\(\Delta \mathrm{Cns}\) (\(\to 0\))
\end{minipage} \\
\midrule\noalign{}
\endhead
\bottomrule\noalign{}
\endlastfoot
Qwen3.5-0.8B (Letters) & Zero-shot & -- & -- & -- & -- & -- & -- \\
Qwen3.5-0.8B (Letters) & Fine-tuned (ISSP) & -- & -- & -- & -- & -- &
-- \\
Kev-0.8B (Pointer) & Off-the-shelf & -- & -- & -- & -- & -- & -- \\
Kev-0.8B (Pointer) & Fine-tuned (ISSP All) & -- & -- & -- & -- & -- &
-- \\
Kev-0.8B (Pointer) & Fine-tuned (ISSP US) & -- & -- & -- & -- & -- &
-- \\
Kev-4B (Pointer) & Off-the-shelf / FT & -- & -- & -- & -- & -- & -- \\
Kev-9B (Pointer) & Off-the-shelf / FT & -- & -- & -- & -- & -- & -- \\
\end{longtable}

\subsubsection{Key Findings (RQ1)}\label{key-findings-rq1}

\begin{itemize}
\tightlist
\item
  \textbf{Off-the-shelf decision models vs.~human distributions:}
  Off-the-shelf decision models are trained for deterministic
  classification/routing and therefore produce overly sharp or
  miscalibrated distributions when asked to represent a heterogeneous
  human subpopulation.
\item
  \textbf{Effect of behavioral fine-tuning:} Fine-tuning with soft
  cross-entropy on empirical human response shares substantially reduces
  both \(D_{\mathrm{KL}}\) and ordinal \(W_1\) distance.
\item
  \textbf{Pointer head vs.~option-letter LM head:} Comparison between
  the pointer-head model and the option-letter LLM on identical
  backbones and training recipes.
\end{itemize}

\subsection{Ordinal Consensus vs.~Human Dispersion
(RQ2)}\label{ordinal-consensus-vs.-human-dispersion-rq2}

\begin{itemize}
\tightlist
\item
  Analysis of \(\mathrm{Cns}(P_{\mathrm{model}})\) versus
  \(\mathrm{Cns}(P_{\mathrm{human}})\) across the 64 Likert
  (\emph{score}) questions in the SubPOP evaluation split.
\item
  Do models capture bimodal human polarization on contentious
  sociopolitical issues, or do they regress toward unimodal
  center-peaked distributions?
\end{itemize}

\subsection{Capturing Inter-Group Demographic Contrasts via
Distributional ANOVA
(RQ3)}\label{capturing-inter-group-demographic-contrasts-via-distributional-anova-rq3}

\begin{itemize}
\tightlist
\item
  \textbf{Persona sensitivity vs.~persona blindness:} Evaluating whether
  models differentiate between contrasting subpopulations (e.g.,
  \emph{Democrat} vs.~\emph{Republican}, \emph{Male} vs.~\emph{Female},
  \emph{Atheist} vs.~\emph{Protestant}, \emph{High} vs.~\emph{Low
  Income}) on questions with large human Wasserstein gaps.
\item
  \textbf{Distributional ANOVA decomposition:} Comparing main effect
  magnitudes (\(\mathbb{E}[E_c]\)), directional consistency
  (\(\mathrm{dPD}\)), and Signal-to-Noise Ratios (\(\mathrm{SNR}\))
  between model predictions and empirical GSS subpopulations.
\item
  Why fixed coarse aggregation (e.g., grouping only by country
  \(\times\) sex \(\times\) age \(\times\) education) causes persona
  insensitivity, and how respondent-level random-subset sampling
  restores demographic conditioning.
\end{itemize}

\subsection{Ablation Studies (RQ4)}\label{ablation-studies-rq4}

\begin{enumerate}
\def\labelenumi{\arabic{enumi}.}
\tightlist
\item
  \textbf{Model Scale and Adaptation Depth:} 0.8B vs.~4B vs.~9B; LoRA
  rank vs.~top-layer unfreezing vs.~full fine-tuning.
\item
  \textbf{Persona Construction and Group Size (\(n_{\min}\)):}
  Single-respondent one-hot targets (\(n=1\)) vs.~aggregated
  subpopulation shares (\(n_{\min} \in \{10, 100\}\)); number of
  conditioning attributes.
\item
  \textbf{Cross-National vs.~Single-Country Training:} Training on all
  310 ISSP country-years vs.~US-only ISSP modules when testing on US GSS
  subpopulations.
\item
  \textbf{Prompt Augmentation vs.~Canonical View:} Impact of shuffling
  option orders and reversing Likert scales during training on test-time
  calibration.
\end{enumerate}

\begin{center}\rule{0.5\linewidth}{0.5pt}\end{center}

\section{Discussion}\label{sec-discussion}

\subsection{When Do Decision Models Make Sense for Sampling Human
Behavior?}\label{when-do-decision-models-make-sense-for-sampling-human-behavior}

\begin{itemize}
\tightlist
\item
  \textbf{Strengths:} Fast, deterministic single-pass inference
  (\(O(1)\) prefill per question); zero syntactic constraint failure
  (\(\mathrm{FailureRate} = 0\) for pointer heads); natural
  representation of subpopulation marginals
  \(P(Y_q \mid \boldsymbol{\lambda})\) when trained with distributional
  cross-entropy.
\item
  \textbf{Fundamental Limitations:}

  \begin{itemize}
  \tightlist
  \item
    \emph{Marginal vs.~Joint Individual Simulation:} Because decision
    models evaluate questions conditionally independently given the
    state
    (\(P(\mathbf{Y} \mid \boldsymbol{\lambda}) = \prod_q P(Y_q \mid \boldsymbol{\lambda})\)),
    sampling from output PMFs reproduces \emph{subpopulation marginals}
    (and convolved construct sums under context isolation) but does not
    automatically capture \emph{within-person cross-question covariance}
    unless conditioning attributes fully explain the correlation.
  \item
    \emph{Out-of-domain attitude generalization:} Limits of inferring
    novel cultural or temporal shifts solely from demographic labels.
  \end{itemize}
\end{itemize}

\subsection{Limitations and Ethical
Considerations}\label{limitations-and-ethical-considerations}

\begin{itemize}
\tightlist
\item
  Risks of over-relying on ``silicon samples'' to replace
  underrepresented human voices in policy or public health research.
\item
  Survey coverage biases in training corpora.
\end{itemize}

\begin{center}\rule{0.5\linewidth}{0.5pt}\end{center}

\section{Conclusion}\label{sec-conclusion}

We investigated whether decision models and constrained next-token LLMs
can serve as calibrated samplers of human behavior. By pairing
respondent-level cross-national survey training (ISSP) with an exact
ordinal evaluation framework---spanning Constraint adherence,
1-Wasserstein distance, ordinal Consensus, and Distributional ANOVA---we
demonstrated both the promise of distributional fine-tuning and the
structural limits of current models in replicating human behavioral
diversity.

\begin{center}\rule{0.5\linewidth}{0.5pt}\end{center}

\section{Appendix}\label{sec-appendix}

\subsection{Proofs and Mathematical Details of Distributional
ANOVA}\label{sec-app-anova}

Let \(\mathcal{D} = \prod_{c \in \mathcal{C}} \{1, \dots, \Lambda_c\}\)
denote the fully crossed factorial design space over demographic and
model factors \(\mathcal{C}\). For a multi-item ordinal instrument
\(\mathbf{Y}_{\boldsymbol{\lambda}} = (Y_{1, \boldsymbol{\lambda}}, \dots, Y_{M, \boldsymbol{\lambda}}) \in \mathcal{Y}^M\),
context isolation between items ensures conditional independence given
condition \(\boldsymbol{\lambda}\) \citep{wadi2025monte, wadi2026every}:
\begin{equation}\protect\phantomsection\label{eq-app-cond-ind}{
P(\mathbf{Y}_{\boldsymbol{\lambda}} = \mathbf{y} \mid \boldsymbol{\lambda}) = \prod_{m=1}^{M} P(Y_{m, \boldsymbol{\lambda}} = y_m \mid \boldsymbol{\lambda}),
}\end{equation} and the exact PMF of the composite construct score
\(S_{\boldsymbol{\lambda}} = \sum_{m=1}^{M} Y_{m, \boldsymbol{\lambda}}\)
is given by discrete convolution
\(P_{S_{\boldsymbol{\lambda}}} = P_{Y_{1, \boldsymbol{\lambda}}} \circledast \cdots \circledast P_{Y_{M, \boldsymbol{\lambda}}}\).
Under the distributional Hoeffding decomposition \citep{wadi2026every},
the expectations of the isolated main-effect distributions
\(E_c(\lambda_c)\) and interaction-effect distributions
\(E_U(\boldsymbol{\lambda}_U)\) recover the classical fixed-effects
ANOVA parameters.

\subsection{Dataset Construction and ISSP Attribute
Harmonization}\label{sec-app-data}

\begin{itemize}
\tightlist
\item
  Full list of 16 ISSP respondent attributes and country-specific
  variable mappings across 2015--2024.
\item
  Overlap analysis between ISSP training items and SubPOP GSS items (29
  overlapping vs.~104 unseen questions).
\end{itemize}


\bibliography{refs/myrefs.bib,refs/custom.bib,refs/custom_2.bib,refs/others.bib,refs/acl-refs.bib,refs/marketing.bib,refs/ai-marketing.bib,refs/silicon-sampling.bib}



\end{document}
